\subsection{Analysis of EMT-Linux Code}
\label{emt:appendix:perf}

\noindent
We analyze the EMT-Linux code to explain 
  the evaluation results of EMT-Linux (e.g., its performance overhead).
We first analyze the source code of EMT-Linux and vanilla Linux.
We use page fault handler
  as the example, as page fault handling 
  contributes to the most kernel activities in our evaluation (Figure~\ref{emt:fig:eval:eval-kern-inst-breakdown}).
Figure~\ref{emt:code:pfh-ve} shows the code snippets of page fault handlers 
  in vanilla Linux 
  and EMT-Linux. 
The two implemenations issues the same set of external function calls
	(both contain one each of \texttt{\small pgd\_offset} and \texttt{\small p\{4,u,m\}d\_alloc})
	and have the same 
    six 64-bit fields on the stack.
Note that EMT enforces translation objects to be passed by pointer rather than by value,
    and thus enables stack allocation and partial updates.
In essence, EMT is an abstraction to existing page table functions
		and do not require an extra ceremony to call.

\begin{figure}[H]
\centering
\begin{subfigure}[b]{.5\textwidth}
\begin{minted}[xleftmargin=17.5pt,linenos,mathescape,breaklines,fontsize=\small,escapeinside=@@]{C}
struct vm_fault { unsigned long address;
  pte_t *pte; pmd_t *pmd; pud_t *pud; pgd_t *pgd; ...
};

static vm_fault_t __handle_mm_fault(...) {
  struct vm_fault vmf = { ...
    .address = address & PAGE_MASK,
  };
  ...
  struct mm_struct *mm = vma->vm_mm;
  ...
  pgd = pgd_offset(mm, address);
  p4d = p4d_alloc(mm, pgd, address);
  vmf.pud = pud_alloc(mm, p4d, address);
  vmf.pmd = pmd_alloc(mm, vmf.pud, address);
  ...
}
\end{minted}
\captionsetup{justification=centering}
\vspace{-20pt}
\caption{Vanilla Linux}
\label{emt:code:pfh-vanilla}
\vspace{-7.5pt}
\end{subfigure}
\begin{subfigure}[b]{.5\textwidth}
\begin{minted}[xleftmargin=17.5pt,linenos,mathescape,breaklines,fontsize=\small,escapeinside=@@]{C}
struct tobj { unsigned long vaddr;
  pte_t *pte; pmd_t *pmd; pud_t *pud; pgd_t *pgd; ...
};

struct vm_fault { ...
  struct tobj tobj;
};

int x86_tdb_find_tobj(struct tdb *tdb, ulong vaddr, struct tobj *out_tobj) {
  ...
  out_tobj->va = vaddr & PAGE_MASK;
  out_tobj->pgd = pgd_offset(tdb, vaddr);
  out_tobj->p4d = p4d_alloc(tdb, out_tobj->pgd, vaddr);
  out_tobj->pud = pud_alloc(tdb, out_tobj->p4d, vaddr);
  out_tobj->pmd = pmd_alloc(tdb, out_tobj->pud, vaddr);
  ...
}

#define tdb_find_tobj x86_tdb_find_tobj

static vm_fault_t __handle_mm_fault(...) {
  struct vm_fault vmf = { ... };
  ...
  if (tdb_find_tobj(mm->tdb, address, &vmf.tobj)) return 1;
  ...
}
\end{minted}
\captionsetup{justification=centering}
\vspace{-20pt}
\caption{EMT-Linux} %\JZc{@siyuan: can you adopt the implementation I put in interface doc?}}
\label{emt:code:pfh-emt}
\end{subfigure}

\vspace{-10pt}
\caption{\bf Simplified code snippets of page fault handlers in (a) vanilla Linux and (b) EMT-Linux.}
\label{emt:code:pfh-ve}
\vspace{-5pt}
\end{figure}

% Although the EMT interface centralizes all page table entries in the \texttt{struct tobj} abstraction
%	and may lead to a multi-word data structure
%	(see Figure~\ref{emt:fig:tobj-x86} for an x86 example),
%		the EMT interface always requires that the \texttt{struct tobj} be passed by pointer rather than by value,
%			thus enabling stack allocation and partial updates.
% By enabling optimizations\footnote{.},
%	the design of the EMT interface allows compilers to generate assemblies with performance approximating vanilla Linux,
%		thus avoiding performance loss.

% To concretize our claim, we use the page fault handler,
%	which generated the most kernel activities in our experiments (Figure \ref{emt:fig:eval:eval-kern-inst-breakdown}),
%	as an example to compare the pre- and post-refactoring code.


\begin{figure}
\centering
\begin{subfigure}[b]{0.23\textwidth}
\centering
\begin{minted}[xleftmargin=0pt,mathescape,breaklines,fontsize=\small,escapeinside=@@]{C}
__handle_mm_fault:
  push rbp
@\diffamv{  mov  rax, rsi}@
  push rbx
@\diffamv{  and  rax, -4096}@
  lea  rbx, [0+rsi*8]
@\diffamv{  mov  rbp, rbx}@
  sub  rsp, 56
@\diffamv{  mov  QWORD PTR [rsp+16], rax}@
  mov  rax, QWORD PTR [rdi]
@\diffoth{}@
@\diffamv{  mov  QWORD PTR [rsp+40], 0}@
  add  rbp, QWORD PTR [rax]
@\diffamv{  mov  QWORD PTR [rsp], rbp}@
@\diffoth{}@
@\diffoth{}@
@\diffoth{}@
@\diffoth{}@
@\diffoth{}@
  je   .L14
  lea  rax, [rbp+0+rbx]
  mov  QWORD PTR [rsp+8], rax
.L6:
  add  rax, rbx
  mov  QWORD PTR [rsp+32], rax
.L8:
  lea  rbp, [rax+rbx]
  mov  QWORD PTR [rsp+24], rbp
  add  rsp, 56
  xor  eax, eax
  pop  rbx
  pop  rbp
  ret
.L14:
  mov  edi, 4096
  call malloc
  test rax, rax
  je   .L3
  mov  QWORD PTR [rsp+8], 0
  mov  edi, 4096
  call malloc
  test rax, rax
  je   .L15
.L5:
  mov  QWORD PTR [rsp+32], 0
.L9:
  mov  edi, 4096
  call malloc
  test rax, rax
  je   .L8
  mov  QWORD PTR [rsp+24], rbp
  add  rsp, 56
  xor  eax, eax
  pop  rbx
  pop  rbp
  ret
.L3:
  mov  QWORD PTR [rsp+8], rbx
  mov  rax, rbx
  test rbx, rbx
  jne  .L6
  mov  edi, 4096
  call malloc
  jmp  .L5
.L15:
  mov  QWORD PTR [rsp+32], rbx
  mov  rax, rbx
  test rbx, rbx
  je   .L9
  jmp  .L8
\end{minted}
\vspace{-10pt}
\caption{vanilla Linux (x86 ASM)}
\label{emt:fig:eval:x86-asm}
\end{subfigure}
\begin{subfigure}[b]{0.23\textwidth}
\centering
\begin{minted}[xleftmargin=0pt,mathescape,breaklines,fontsize=\small,escapeinside=@@]{C}
__handle_mm_fault:
  push rbp
@\diffoth{}@
  push rbx
@\diffoth{}@
  lea  rbx, [0+rsi*8]
@\diffoth{}@
  sub  rsp, 56
@\diffoth{}@
  mov  rax, QWORD PTR [rdi]
@\diffdmv{  mov  QWORD PTR [rsp+8], 0}@
@\diffdel{  mov  rax, QWORD PTR [rax]}@
  mov  rbp, QWORD PTR [rax]
@\diffoth{}@
@\diffdmv{  mov  rax, rsi}@
@\diffdmv{  and  rax, -4096}@
@\diffdmv{  add  rbp, rbx}@
@\diffdmv{  mov  QWORD PTR [rsp], rax}@
@\diffdmv{  mov  QWORD PTR [rsp+40], rbp}@
  je   .L16
  lea  rax, [rbp+0+rbx]
  mov  QWORD PTR [rsp+32], rax
.L6:
  add  rax, rbx
  mov  QWORD PTR [rsp+24], rax
.L8:
  lea  rbp, [rax+rbx]
  mov  QWORD PTR [rsp+16], rbp
  add  rsp, 56
  xor  eax, eax
  pop  rbx
  pop  rbp
  ret
.L16:
  mov  edi, 4096
  call malloc
  test rax, rax
  je   .L3
  mov  QWORD PTR [rsp+32], 0
  mov  edi, 4096
  call malloc
  test rax, rax
  je   .L17
.L5:
  mov  QWORD PTR [rsp+24], 0
.L9:
  mov  edi, 4096
  call malloc
  test rax, rax
  je   .L8
  mov  QWORD PTR [rsp+16], rbp
  add  rsp, 56
  xor  eax, eax
  pop  rbx
  pop  rbp
  ret
.L3:
  mov  QWORD PTR [rsp+32], rbx
  mov  rax, rbx
  test rbx, rbx
  jne  .L6
  mov  edi, 4096
  call malloc
  jmp  .L5
.L17:
  mov  QWORD PTR [rsp+24], rbx
  mov  rax, rbx
  test rbx, rbx
  je   .L9
  jmp  .L8
\end{minted}
\vspace{-10pt}
\caption{EMT-Linux (x86 ASM)}
\label{emt:fig:eval:emt-asm}
\end{subfigure}
\vspace{-10pt}
\caption{{\bf Compiled code of page fault handler snippets 
  of vanilla Linux and EMT-Linux.} 
  Main difference is highlighted in dark color; moved lines are in light color.
  For demonstration, we used very simplified code snippets, which does not reflect 
  the complexity of real page fault handlers.}
\label{emt:fig:eval:asm}
\end{figure}

Figure~\ref{emt:fig:eval:asm} shows 
  the x86 assembly generated by compiling the code snippets in Figure~\ref{emt:code:pfh-ve}.
We use GCC v14.2 with -O2 optimization as Linux requires the compiler optimization level 
  to be at least -O2 since version 0.99.14o~\cite{linux:require-o2}.
% Figure~\ref{emt:fig:eval:x86-asm} and \ref{emt:fig:eval:emt-asm} shows the compiled code in x86 \jzx{ASM}{assembly}
%  of \jz{our page fault handler snippets ()} \jzx{in}{with} vanilla Linux~\cite{linux:code:__handle_mm_fault:x86} and in EMT-Linux~\cite{linux:code:__handle_mm_fault:emt} \jz{APIs}, respectively.
We can see in Figure~\ref{emt:fig:eval:asm} that
  EMT-Linux only adds one extra memory read instructions
  due to the dereferncing of \texttt{\small mm->tdb} which is a pointer to the translation database instance
  (in vanilla Linux, the kernel uses \texttt{mm->pgd} which simply store the root address of radix-tree page table).
\schai{We chose to use a pointer for \texttt{\small tdb}instead of inlining the struct, because the translation database is defined by the MMU driver and may have different size, scope, and structure.}
Other changes are merely instruction or field location reorderings.

% We further examine the compilation results of the above-shown code
%	and the results show that the example code in the two figures produced nearly identical compilation results when using GCC 14.2 and -O2 optimization (see \cite{linux:code:__handle_mm_fault:x86} and \cite{linux:code:__handle_mm_fault:emt}).

We count the instructions of the page fault handlers
	in the real EMT-Linux and vanilla Linux artifacts ({\it not} simplified ones as shown in Figure~\ref{emt:code:pfh-ve}).
The page fault handler core routine (\texttt{\small \_\_handle\_mm\_fault}) in vanilla Linux
	has 1165 instructions, including 274 memory-load instructions and 64 memory-store instructions.
The \texttt{\small \_\_handle\_mm\_fault} routine in EMT-Linux has 1172 instructions, including 281 memory-load instructions 
  and 61 memory-store instructions.

We observe that EMT-Linux and vanilla Linux have similar cache efficiency.
For example, on the Memcached workload, 
  vanilla Linux and EMT-Linux have cache hit rates of 51.3\% and 51.7\%,
  with L1I cache hit rates of 90.8\% and 91.2\%, respectively.
We think that the small EMT interface, which centralizes frequently called 
  MMU driver code,
  is cache friendly.
%  with a small interface, % \jz{and centralize page-table related data into dedicated data structures such as \texttt{struct tobj}},
% EMT-Linux indirectly reduces the kernel's cache miss rate
% by creating hot paths. % compared to vanilla Linux.
% Fundamentally, OS code, due to its relatively short running time,
%	has less hot caches~\cite{linux:code:icache}, compared to user programs.
% Since EMT's basic function API contains only 15 functions that are frequently invoked
%  and encapsulates translation-related data, 
%	EMT-Linux is friendly to instructions and data caches.
% \jz{On the other hand, EMT uses compact data sturctures such as \jz{struct tobj} for page table operations,
%	improving the locality of accessing such data.
% This locality also helps reduce the chance of cache misses.}
% \JZx{This concentration of calls indirectly reduces the miss rate of caches.}
% We find that EMT-Linux has slightly lower cache miss rates (measured by the Linux Perf tool~\cite{perf}).
% For memcached workload, we evaluated the kernel's cache miss rate by using the .
% The reduction in cache miss rate offsets the slight increase in instruction count caused 
%  by using the EMT interface.

Lastly, EMT employs a macro-based function redirection interface for
  transparent function rewriting, enabling MMU drivers to define macros or inline functions when necessary.
With this feature, EMT can reduce the overhead of calling small functions.
In our implementation,
	we strictly follow the Linux coding style~\cite{linux:code:icache} and only use function inlining
	for functions with no more than three lines
	or functions that can perform compile-time code elimination.

\begin{comment}

\begin{figure}[H]
\begin{minted}[xleftmargin=0pt,mathescape,breaklines,fontsize=\small,escapeinside=@@]{C}
// In MMU driver implementation file
int x86_tdb_find(...) {
  // ...
}

// In MMU driver header file
#define tobj_equal(tobj1, tobj2) \
  (((tobj1) == (tobj2)) || \
    !memcmp((tobj1), (tobj2), sizeof(struct tobj)))

static inline x86_tobj_read_attr(...) {
  // ...
}
#define tobj_read_attr x86_tobj_read_attr

int x86_tdb_find(...);
#define tdb_find x86_tdb_find

// In EMT default implementation file
#ifndef tobj_equal
#error Basic function tobj_equal is not defined
#endif

#ifndef tobj_read_attr
#error Basic function tobj_equal is not defined
#endif

#ifndef tdb_find
#error Basic function tdb_find is not defined
#endif

\end{minted}
\vspace{-12.5pt}
\caption{{\bf Implementation of EMT page fault handling routine.}}
\label{emt:fig:emt-redir}
\vspace{-5pt}
\end{figure}

\end{comment}